\documentclass{article}
\usepackage{graphicx} % Required for inserting images

%https://www.sarmate.net/demo/mathpad_tex_edit.php?lang=en

\title{Math1800 Group Project}
\author{Tyson Wheeler}
\date{Semester 2, 2026}

\begin{document}

\maketitle

\section{Introduction}
In forensic investigations, estimating the time of death can be essential for police investigation. This model examines how the temperature of a deceased body changes over time as heat is exchanged with the surrounding environment. Focusing particularly on when environmental conditions change. The aim of this report is to improve the basic estimation of the postmortem interval by accounting for environmental conditions like the change in the ambient temperature throughout the day.
Newton’s law of cooling assumes that the surface temperature of an object changes at a rate proportional to the difference between the object and the ambient temperature. However, this relies on the constant conditions, which are rarely present in real life situations, reducing the precision of this model in estimating the time of death.
Algor mortis is known as the cooling of the body after death. The rate at which this occurs is affected by environmental factors, particularly ambient temperature, which can cool the body off at differing rates. Since Newton’s law assumes relatively constant environmental conditions and a simplified relationship between temperature and heat loss, it cannot fully account for the complexities of a real body.
This model explores how a deceased body may initially cool and subsequently warm if environmental as the ambient temperature changes due to the relationship between time, ambient temperature, and body temperature.
An accurate estimation of time of death can provide police with essential information such as establishing a timeline, verifying alibis, finding, and identifying possible suspects. Thus, the aim of this model is to improve the accuracy of predicting the postmortem interval by considering external factors.

\section{Mathematical Formulation and Solutions of Equations}

\section{Discussion of Results}

\section{Conclusion}

\section{References}

\end{document}

